\begin{table}[h]
  \centering
  \caption{면접 기반 선행 연구와 본 SNS 자료의 분포 비교. 본 연구의 과정 코드는 문보경$\cdot$장성숙(2001)의 범주를 출발점으로 만든 것이며, 두 자료원의 결과가 같은 모집단의 일치를 의미하지는 않는다.}
  \label{tab:compare}
  \small
  \begin{tabularx}{\columnwidth}{@{}>{\raggedright\arraybackslash}p{0.16\columnwidth} >{\raggedright\arraybackslash}X >{\raggedright\arraybackslash}X@{}}
    \toprule
    \textbf{측면} & \textbf{면접 자료 (문보경$\cdot$장성숙, 2001)} & \textbf{본 SNS 자료} \\
    \midrule
    자료원 & 내담자 면접$\cdot$개방형 질문지 & 한국어 공개 X 경험 글 \\
    분석 단위 & 사례 단위 사건 회고 & 게시물 단위 경험 글 \\
    부정 표지 & 원치 않는 반응, 해결책 부재, 상담자의 요구, 상처 경험 & 강제성 0.350, 무시당함 0.470, 상처 0.231 (부정 117건 기준) \\
    분석 깊이 & 사례별 깊이 있는 합의 부호화 & 분포 수준의 탐색적 부호화 \\
    \bottomrule
  \end{tabularx}
\end{table}
